\documentclass[11pt,a4paper]{article}

\usepackage[margin=2.5cm]{geometry}
\usepackage{graphicx}
\usepackage{amsmath}
\usepackage{hyperref}
\usepackage{booktabs}
\usepackage{enumitem}
\usepackage{titlesec}
\usepackage{caption}
\usepackage{subcaption}
\usepackage{float}
\usepackage{natbib}
\usepackage{listings}

\titleformat{\section}{\large\bfseries}{\thesection}{1em}{}
\titleformat{\subsection}{\normalsize\bfseries}{\thesubsection}{1em}{}

\title{Métodos Basados en el Conocimiento Aplicado a Empresas\\
       Módulo de Metaheurística}

\author{\\\\
\textbf{Integrantes:}\\\\
Andrés Augusto Aragón Dimas - UO316049 \\
Alberto Martínez Olivar - UO282069\\ 
Manuel De La Uz González - UO288991\\ 
Juan Camilo España Lopera - U0318200\\\\\\
\textbf{Profesor:}\\\\
Jorge Puente Peinador\\\\\\
Universidad de Oviedo\\
}

\date{\today}

\begin{document}

\maketitle

\newpage
\tableofcontents
% -------------------------------------------------\section{Introducción y Contexto del Problema}

\newpage
\section{Introducción y Contexto del Problema}
\section{Modelamiento Formal del Problema de Drones}

\subsection*{Índices}

\begin{alignat}{2}
&i \in I \qquad && \text{índice de drones} \\
&j \in J \qquad && \text{índice de puntos: barridos y estación de control} \\
&k \in K \qquad && \text{índice del orden de visita de puntos en una ruta}
\end{alignat}

\noindent siendo $j = 1$ la estación de control y los índices restantes los extremos de cada barrido.

\subsection*{Parámetros}

\begin{alignat}{2}
&n = |I| \qquad && \text{número de drones disponibles} \\
&m = |J| = 2 \times |\text{barridos}| + 1 \qquad && \text{número total de puntos incluyendo la estación de control} \\
&D_{j_1 j_2} \qquad && \text{distancia euclídea entre el punto } j_1 \text{ y el punto } j_2, \quad j_1, j_2 \in J
\end{alignat}

\subsection*{Variable de Decisión}

\begin{equation}
    X_{ijk} = \begin{cases}
        1 & \text{si el dron } i \text{ visita el punto } j \text{ en el orden } k \\
        0 & \text{en caso contrario}
    \end{cases}
\end{equation}

\subsection*{Variables Auxiliares}

\begin{equation}
    D_i = \sum_{j_1 \in J} \sum_{j_2 \in J} \sum_{k \in K} X_{ij_1k} \cdot X_{ij_2,k+1} \cdot D_{j_1 j_2} \qquad \forall i \in I
\end{equation}

\noindent donde $D_i$ representa la distancia total recorrida por el dron $i$.

\bigskip
\noindent\textbf{Nota sobre no linealidad:} la expresión de $D_i$ contiene productos de variables binarias, lo que introduce no linealidad en el modelo. Esta formulación se mantiene de forma descriptiva con el único propósito de clarificar la estructura del problema; la resolución se realiza mediante metaheurísticas, que no requieren linealidad en el modelo matemático.


\section{Modelamiento del Algoritmo Genético}

\section{Descripción del Código Implementado}

\section{Fase 1: Evaluación de la Convergencia Práctica}

\section{Fase 2: Análisis Paramétrico}

\subsection{Factores Mínimos de Estudio}

\section{Indicadores y Análisis de Resultados}

\section{Contraste con Búsqueda Aleatoria}

\section{Interpretación de Soluciones Representativas del Frente de Pareto}

\section{Conclusiones}

\section{Aportaciones Individuales}

\section{Referencias y Recursos Utilizados}

Ejemplo citación \cite{chatterjee1999recursive}


\bibliographystyle{plainnat} 
\bibliography{ref}

\end{document}